\documentclass[11pt]{article}

\usepackage[margin=1in]{geometry}
\usepackage{amsmath,amssymb,amsthm}
\usepackage{graphicx}
\usepackage[hidelinks]{hyperref}

\newtheorem{theorem}{Theorem}
\newtheorem{corollary}{Corollary}

\newcommand{\R}{\mathbb{R}}
\newcommand{\Id}{\mathrm{Id}}
\newcommand{\prox}{\operatorname{prox}}
\newcommand{\argmin}{\operatorname*{arg\,min}}
\newcommand{\E}{\mathbb{E}}

\title{A Non-Proximal Firmly Nonexpansive Expected-Risk Minimizer}
\author{Run \texttt{fa\_banach\_001}, agent \texttt{agent\_lane\_04}}
\date{2026-07-19}

\begin{document}
\maketitle

\section*{Status}

This packet gives a candidate counterexample, likely valid, to the universal
form of a future-work question in Bredies--Chirinos-Rodriguez--Naldi
\cite{BCN}.  It does not contradict their numerical experiments.  It shows that
extra statistical or structural assumptions are needed if one wants a theorem
forcing learned firmly nonexpansive maps to be proximity operators.

\section*{Source Question}

The source paper studies the expected-risk problem
\[
        N^* \in \argmin_{N\in\mathcal N}
        F(N), \qquad
        F(N)=\int_{\mathcal X\times \mathcal X}
             \|N(\bar x)-\bar y\|^2\,d\mu(\bar x,\bar y),
\]
where $\mathcal N$ is the class of nonexpansive maps.  The corresponding firmly
nonexpansive learned operator is obtained by the change of variables
$T=(\Id+N)/2$.

Near the end of the ``Third experiment: proximity operators'' discussion, the
authors ask whether learned firmly nonexpansive operators preferentially
converge to proximity operators, or, more strongly, whether proximity operators
are the only objects in the argmin of this expected-risk problem for ``real''
data distributions $\mu$.

\begin{center}
\includegraphics[width=\linewidth]{figures/open_problem_crop.png}
\end{center}

\section*{Proof Intuition}

The obstruction is rotational.  A proximity operator of a convex function is a
conservative gradient field, hence its circulation around every closed curve is
zero.  But the average of the identity with a quarter-turn rotation is still
firmly nonexpansive while retaining nonzero circulation.  If the data are
generated exactly from that map and the input distribution has full support,
zero expected risk forces every minimizer to agree with the rotational map
everywhere.  Thus the unique learned firmly nonexpansive operator is not
proximal.

\section*{Counterexample}

Let
\[
        R=\begin{pmatrix}0&-1\\ 1&0\end{pmatrix},
        \qquad
        T=\frac{\Id+R}{2}.
\]
Let $X$ be a standard Gaussian random vector in $\R^2$ and set $Z=T X$.
Equivalently, in the paper's nonexpansive formulation, set
\[
        N=2T-\Id=R,\qquad Y=NX=RX.
\]
Let $\mu'$ be the law of $(X,Z)$ and let $\mu$ be the law of $(X,Y)$.
Both distributions have finite second moments and therefore satisfy the source
paper's expected-risk integrability hypothesis.

\begin{theorem}
For the direct firmly nonexpansive risk
\[
        F'(S)=\E\|S(X)-Z\|^2,
        \qquad S \text{ firmly nonexpansive},
\]
the unique minimizer is $S=T$.  The minimizer $T$ is not a proximity operator of
any proper lower semicontinuous convex function on $\R^2$.
\end{theorem}

\begin{proof}
First, $T$ is firmly nonexpansive.  For every $h\in\R^2$,
\[
        \|Th\|^2
        =\frac14\|h+Rh\|^2
        =\frac12\|h\|^2
        =\langle Th,h\rangle,
\]
because $R$ is orthogonal and skew-symmetric.  This is exactly the linear
firm-nonexpansiveness condition
$\|Th\|^2\leq \langle Th,h\rangle$.

Since $Z=TX$, this $T$ has risk $F'(T)=0$, so the minimum risk is $0$.  If $S$
is another firmly nonexpansive minimizer, then
$\E\|S(X)-TX\|^2=0$, hence $S(X)=TX$ almost surely.  Firmly nonexpansive maps
are continuous.  The standard Gaussian has positive mass on every nonempty open
ball in $\R^2$.  Therefore, if $Sx\neq Tx$ at any point $x$, continuity would
give a nonempty open ball on which $S$ and $T$ differ, contradicting almost-sure
equality.  Hence $S=T$ on all of $\R^2$.

It remains to show that $T$ is not proximal.  A proximity operator of a proper
lower semicontinuous convex function is the gradient of a convex $C^1$ potential
\cite{BCN,BC}; in particular, its circulation around every closed $C^1$ curve is
zero.  On the unit circle $\gamma(t)=(\cos t,\sin t)$, $0\leq t\leq 2\pi$, we
have $\gamma'(t)=R\gamma(t)$ and
\[
        \int_\gamma T(x)\cdot dx
        =
        \int_0^{2\pi}
        \left\langle \frac{\gamma(t)+R\gamma(t)}2,
        R\gamma(t)\right\rangle\,dt
        =
        \int_0^{2\pi}\frac12\,dt
        =\pi.
\]
The circulation is nonzero, so $T$ is not conservative and hence cannot be a
proximity operator.
\end{proof}

\begin{corollary}
In the paper's nonexpansive formulation (CP), the distribution $\mu=\mathcal
L(X,RX)$ has the unique minimizer $N=R$.  The associated learned firmly
nonexpansive operator $(\Id+N)/2$ is non-proximal.
\end{corollary}

\begin{proof}
The rotation $R$ is nonexpansive and has zero risk for (CP).  The same
full-support argument used above shows that any nonexpansive zero-risk minimizer
must agree with $R$ everywhere.  The associated firmly nonexpansive map is
$(\Id+R)/2=T$, which is non-proximal by the theorem.
\end{proof}

\section*{Verification Notes}

The proof is analytic.  The included script
\texttt{code/verify\_linear\_counterexample.py} checks:
\begin{itemize}
    \item the matrix inequality $T^*T\leq (T+T^*)/2$;
    \item $\|2T-\Id\|=1$, so the corresponding (CP) variable is nonexpansive;
    \item the unit-circle circulation is numerically $\pi$.
\end{itemize}

Command used:
\[
\texttt{conda run --no-capture-output -n sandbox python code/verify\_linear\_counterexample.py}.
\]

\section*{Novelty Check and Limitations}

Local indexes searched: \texttt{registry\_index.tsv}, \texttt{solutions/index.tsv},
\texttt{attempts/index.tsv}, and \texttt{proof\_gaps/index.tsv}, using
\texttt{2407.14156}, \texttt{Learning Firmly}, \texttt{proximal argmin},
\texttt{nonproximal}, and close firmly-nonexpansive/proximity phrases.  No
duplicate packet or attempt was found.

External web search on 2026-07-19 for the exact phrase ``are proximity
operators the only objects in the argmin'', together with close variants
involving ``Learning Firmly Nonexpansive Operators'', ``firmly nonexpansive'',
``proximity operators'', and ``expected risk'', found the source paper and
general plug-and-play/proximity-operator references, but no exact answer to this
future-work question.

The result refutes only the universal form of the question.  It does not rule
out the softer ``preferential convergence'' possibility, nor positive theorems
for natural-image distributions, separable gradient data, symmetry assumptions,
or additional constraints that enforce conservativity/cyclic monotonicity.

\section*{Human Review Recommendation}

Reviewers should check that the source question is interpreted as a universal
mathematical claim under the expected-risk framework, and that the distinction
between the source paper's nonexpansive variable $N$ and the associated firmly
nonexpansive variable $T=(\Id+N)/2$ is acceptable.  The proof itself reduces to
standard two-dimensional Hilbert-space calculations.

\begin{thebibliography}{9}

\bibitem{BCN}
K. Bredies, J. Chirinos-Rodriguez, and E. Naldi,
\emph{Learning Firmly Nonexpansive Operators},
arXiv:2407.14156v3, 2025.

\bibitem{BC}
H. H. Bauschke and P. L. Combettes,
\emph{Convex Analysis and Monotone Operator Theory in Hilbert Spaces},
2nd ed., Springer, 2017.

\end{thebibliography}

\end{document}
